\clearpage
\begin{explanationbox}
\textbf{\textcolor{CExplanation}{先看六个常见奇函数：对称相同，零处是否有定义不同}}

\begin{center}\begin{tikzpicture}[font=\small,>=stealth,line join=round]
\draw[gray!40,rounded corners=3pt] (0,0) rectangle (14.2,8.4);\node at (7.1,7.800000000000001) {常见奇函数：先看关于原点的对称，再看能否代入零};
\node at (2.3,6.9) {$y=x$};\begin{scope}[shift={(2.3,4.3)},x=1cm,y=1cm]\draw[->,gray] (-1.55,0)--(1.55,0) node[right]{$x$};\draw[->,gray] (0,-1.75)--(0,1.75) node[above]{$y$};\draw[blue,thick,domain=-1.4:1.4,samples=150] plot(\x,{\x});\fill[orange!85!black] (0,0) circle(2.2pt);\end{scope}\node[text width=4.0cm,align=center] at (2.3,1.7) {$D=\mathbb R$，$f(0)=0$};\node at (6.95,6.9) {$y=x^3$};\begin{scope}[shift={(6.95,4.3)},x=1cm,y=0.85cm]\draw[->,gray] (-1.55,0)--(1.55,0) node[right]{$x$};\draw[->,gray] (0,-2.058823529411765)--(0,2.058823529411765) node[above]{$y$};\draw[blue,thick,domain=-1.2:1.2,samples=150] plot(\x,{(\x)^3});\fill[orange!85!black] (0,0) circle(2.2pt);\end{scope}\node[text width=4.0cm,align=center] at (6.95,1.7) {$D=\mathbb R$，$f(0)=0$};\node at (11.600000000000001,6.9) {$y=\sqrt[3]{x}$};\begin{scope}[shift={(11.600000000000001,4.3)},x=1cm,y=1.25cm]\draw[->,gray] (-1.55,0)--(1.55,0) node[right]{$x$};\draw[->,gray] (0,-1.4)--(0,1.4) node[above]{$y$};\draw[blue,thick,domain=-1.4:1.4,samples=150] plot(\x,{sign(\x)*pow(abs(\x),1/3)});\fill[orange!85!black] (0,0) circle(2.2pt);\end{scope}\node[text width=4.0cm,align=center] at (11.600000000000001,1.7) {$D=\mathbb R$，$f(0)=0$};
\node[text width=12.8cm,align=center] at (7.1,.65) {每个小图独立取合适的纵轴尺度，不能据此比较陡缓；橙点均为原点。};
\end{tikzpicture}\end{center}


\begin{center}\begin{tikzpicture}[font=\small,>=stealth,line join=round]
\draw[gray!40,rounded corners=3pt] (0,0) rectangle (14.2,8.4);\node at (7.1,7.800000000000001) {常见奇函数：先看关于原点的对称，再看能否代入零};
\node at (2.3,6.9) {$y=\sin x$};\begin{scope}[shift={(2.3,4.3)},x=1cm,y=1.5cm]\draw[->,gray] (-1.55,0)--(1.55,0) node[right]{$x$};\draw[->,gray] (0,-1.1666666666666667)--(0,1.1666666666666667) node[above]{$y$};\draw[blue,thick,domain=-1.45:1.45,samples=150] plot(\x,{sin(deg(\x))});\fill[orange!85!black] (0,0) circle(2.2pt);\end{scope}\node[text width=4.0cm,align=center] at (2.3,1.7) {$D=\mathbb R$，$f(0)=0$};\node at (6.95,6.9) {$y=1/x$};\begin{scope}[shift={(6.95,4.3)},x=1cm,y=1cm]\draw[->,gray] (-1.55,0)--(1.55,0) node[right]{$x$};\draw[->,gray] (0,-1.75)--(0,1.75) node[above]{$y$};\draw[blue,thick,domain=-1.4:-0.6,samples=150] plot(\x,{1/\x});\draw[blue,thick,domain=0.6:1.4,samples=150] plot(\x,{1/\x});\end{scope}\node[text width=4.0cm,align=center] at (6.95,1.7) {$D=\mathbb R\setminus\{0\}$\\$f(0)$ 不存在};\node at (11.600000000000001,6.9) {$y=\frac{x}{1+x^2}$};\begin{scope}[shift={(11.600000000000001,4.3)},x=1cm,y=3cm]\draw[->,gray] (-1.55,0)--(1.55,0) node[right]{$x$};\draw[->,gray] (0,-0.5833333333333334)--(0,0.5833333333333334) node[above]{$y$};\draw[blue,thick,domain=-1.4:1.4,samples=150] plot(\x,{\x/(1+(\x)^2)});\fill[orange!85!black] (0,0) circle(2.2pt);\end{scope}\node[text width=4.0cm,align=center] at (11.600000000000001,1.7) {$D=\mathbb R$，$f(0)=0$};
\node[text width=12.8cm,align=center] at (7.1,.65) {每个小图独立取合适的纵轴尺度；$1/x$ 的两支不相连，原点不在图像上。};
\end{tikzpicture}\end{center}


\end{explanationbox}
\clearpage
\begin{explanationbox}
  \textbf{\textcolor{CExplanation}{一、这个结论为什么会出现？}}

  我们知道，奇函数有一条规则：把横坐标 $x$ 换成 $-x$，函数值也要变成原来的相反数。
  例如，如果 $f(2)=5$，那么 $f(-2)=-5$。
  所以奇函数的图像，总是围绕原点成对出现。

  可是 $x=0$ 与别的位置不同，因为
  \[
    -0=0.
  \]
  把 $x=0$ 变成它的相反数，横坐标竟然没有变化；
  但奇函数的规则还要求纵坐标变成相反数。
  \textbf{同一个横坐标只能有一个函数值}，
  于是这个函数值必须等于自己的相反数。
  \[
    f(0)=-f(0)\quad\Longrightarrow\quad f(0)=0.
  \]
  \textbf{原来并不是要死记 $f(0)=0$；而是 $0$ 恰好是“变成相反数后仍然不变”的那个横坐标。}

  \medskip
  \textbf{\textcolor{CExplanation}{二、用图像想象这件事}}
\begin{center}\begin{tikzpicture}[font=\small,>=stealth,line join=round]
\draw[gray!40,rounded corners=3pt] (0,0) rectangle (14.2,8.5);\node at (7.1,7.9) {在同一条竖线上，对称点必须重合};
\node at (3.5,6.85){假设 $c\ne0$：同一个输入出现两个值};\node at (10.6,6.85){$c=0$：两个点合为原点};\begin{scope}[shift={(3.5,3.8)},x=1cm,y=1cm]\draw[->,gray] (-2.1,0)--(2.1,0) node[right]{$x$};\draw[->,gray] (0,-2.0)--(0,2.0) node[above]{$y$};\fill[red] (0,1.3) circle(2.2pt) node[right=5pt] {$(0,c)$};\fill[red] (0,-1.3) circle(2.2pt) node[right=5pt] {$(0,-c)$};\draw[red,dashed] (0,-1.3)--(0,1.3);\end{scope}\begin{scope}[shift={(10.5,3.8)},x=1cm,y=1cm]\draw[->,gray] (-2.1,0)--(2.1,0) node[right]{$x$};\draw[->,gray] (0,-2)--(0,2) node[above]{$y$};\draw[blue,thick,domain=-1.2:1.2,samples=150] plot(\x,{(\x)^3});\fill[orange!85!black] (0,0) circle(2.2pt) node[below right=5pt] {$(0,0)$};\end{scope}\node at (3.5,1.7){左侧不能作为函数的图像};\node at (10.6,1.7){一个横坐标只对应一个纵坐标};
\node[text width=12.8cm,align=center] at (7.1,.65) {左图为 $c\ne0$ 的假设示意，并非函数图像；$c<0$ 时上下两点交换。};
\end{tikzpicture}\end{center}



  看左图：若 $c\ne0$，原点对称就要求在同一条竖线上同时出现 $(0,c)$ 与 $(0,-c)$；
  这违反一个输入只能对应一个函数值的要求。
  看右图：只有两点重合才可行，因此
  \[
    c=-c,\qquad c=0.
  \]
  因此，奇函数的图像\textbf{只要与 $y$ 轴有交点，这个交点必然就是原点 $(0,0)$}。

  \medskip
  
\end{explanationbox}
\clearpage
\begin{explanationbox}
\textbf{\textcolor{CExplanation}{三、它在解题时有什么用？}}
\begin{center}\begin{tikzpicture}[font=\small,>=stealth,line join=round]
\draw[gray!40,rounded corners=3pt] (0,0) rectangle (14.2,8.5);\node at (7.1,7.9) {非零截距会怎样破坏原点对称？};
\begin{scope}[shift={(4.0,3.8)},x=1.2cm,y=0.9cm]\draw[->,gray] (-1.9,0)--(1.9,0) node[right]{$x$};\draw[->,gray] (0,-2.8)--(0,2.8) node[above]{$y$};\draw[blue,thick,domain=-1.05:1.05,samples=150] plot(\x,{(\x)^3+1});\draw[violet,dashed,domain=-1.05:1.05,samples=100] plot(\x,{(\x)^3-1});\fill[blue] (0,1) circle(2.2pt) node[above left=5pt] {$(0,1)$};\draw[violet,fill=white] (0,-1) circle(2.5pt);\node[below right=8pt,violet] at (0,-1){$(0,-1)$};\end{scope}\node[align=left,anchor=west] at (8.0,4.1){蓝实线：$y=x^3+1$\\[11pt]紫虚线：旋转半圈后的图像\\[11pt]两图不重合，不能是奇函数\\[11pt]只看 $f(0)=1\ne0$ 就能排除};
\node[text width=12.8cm,align=center] at (7.1,.65) {虚线为蓝图的原点对称像 $y=x^3-1$，不是原函数的另一部分。};
\end{tikzpicture}\end{center}


  \begin{description}[style=nextline, leftmargin=2.8em, labelsep=0.5em, itemsep=0.45em]
    \item[\textbf{用途一：不求完整表达式，直接求值。}]
      题目说 $f$ 是奇函数、定义域为全体实数，那么立即写出 $f(0)=0$。
    \item[\textbf{用途二：求未知参数。}]
      如果 $f(0)=a-3$，又知道 $f$ 是奇函数且在 $0$ 处有定义，就能立刻得到 $a=3$。
    \item[\textbf{用途三：快速排除。}]
      若能算出 $f(0)=4$，就可以断定这个函数不是奇函数，不必再代入一长串 $f(-x)$。
  \end{description}

  \medskip
  
\end{explanationbox}
\clearpage
\begin{explanationbox}
\textbf{\textcolor{CExplanation}{四、联系之前学过的知识}}
\begin{center}\begin{tikzpicture}[font=\small,>=stealth,line join=round]
\draw[gray!40,rounded corners=3pt] (0,0) rectangle (14.2,8.5);\node at (7.1,7.9) {易错五：偶函数在零处的值，可以不是零};
\node at (3.5,6.85){奇函数 $y=x^3$};\node at (10.6,6.85){偶函数 $y=x^2+2$};\begin{scope}[shift={(3.5,3.1)},x=0.95cm,y=0.75cm]\draw[->,gray] (-2.2,0)--(2.2,0) node[right]{$x$};\draw[->,gray] (0,-1.5)--(0,4.1) node[above]{$y$};\draw[blue,thick,domain=-1.1:1.55,samples=150] plot(\x,{(\x)^3});\fill[orange!85!black] (0,0) circle(2.2pt) node[below right=5pt] {$(0,0)$};\end{scope}\begin{scope}[shift={(10.5,3.1)},x=0.95cm,y=0.75cm]\draw[->,gray] (-2.2,0)--(2.2,0) node[right]{$x$};\draw[->,gray] (0,-1.5)--(0,4.1) node[above]{$y$};\draw[blue,thick,domain=-1.4:1.4,samples=150] plot(\x,{(\x)^2+2});\fill[orange!85!black] (0,2) circle(2.2pt) node[below right=5pt] {$(0,2)$};\fill[orange!85!black] (-1,3) circle(2.2pt);\fill[orange!85!black] (1,3) circle(2.2pt);\draw[gray,dashed] (-1,3)--(1,3);\end{scope}
\node[text width=12.8cm,align=center] at (7.1,.65) {两图同尺度；左图要求相反数配对，右图只要求左右高度相同。};
\end{tikzpicture}\end{center}



  \textbf{（1）与偶函数对照。}
  偶函数满足 $f(-x)=f(x)$。代入 $x=0$，只会得到 $f(0)=f(0)$；
  这句话无论函数值是多少都成立，所以\textbf{偶函数在 $0$ 处的函数值不必为 $0$}。
  例如 $f(x)=x^2+2$ 是偶函数，$f(0)=2$。

  
\end{explanationbox}
\clearpage
\begin{explanationbox}
\textbf{\textcolor{CExplanation}{四、联系之前学过的知识（续）：零点的配对}}

\begin{center}\begin{tikzpicture}[font=\small,>=stealth,line join=round]
\draw[gray!40,rounded corners=3pt] (0,0) rectangle (14.2,8.5);\node at (7.1,7.9) {易错六：原点必为零点，不代表它是唯一零点};
\begin{scope}[shift={(4.2,3.8)},x=1.8cm,y=1.5cm]\draw[->,gray] (-1.9,0)--(1.9,0) node[right]{$x$};\draw[->,gray] (0,-1.7)--(0,1.7) node[above]{$y$};\draw[blue,thick,domain=-1.4:1.4,samples=150] plot(\x,{(\x)^3-\x});\fill[orange!85!black] (-1,0) circle(2.2pt) node[below=5pt] {$-1$};\fill[orange!85!black] (0,0) circle(2.2pt) node[above right=8pt] {$0$};\fill[orange!85!black] (1,0) circle(2.2pt) node[below=5pt] {$1$};\draw[gray,dotted] (-1.2,-1.15)--(1.2,-1.15);\draw[<->,orange!85!black] (-1,-1.05)--(1,-1.05);\end{scope}\node[align=left,anchor=west] at (8.4,4.1){$f(x)=x^3-x$\\[11pt]$=x(x-1)(x+1)$\\[11pt]零点：$-1,0,1$\\[11pt]两侧零点成对};
\node[text width=12.8cm,align=center] at (7.1,.65) {曲线直接穿过三个标出的零点；$f(-x)=-f(x)$。};
\end{tikzpicture}\end{center}

\textbf{（2）与奇函数的零点对照。}
  若 $a$ 和 $-a$ 在定义域内，且奇函数满足 $f(a)=0$，那么
  \[
    f(-a)=-f(a)=0.
  \]
  图中 $-1$ 与 $1$ 是一对非零零点，而 $0$ 是自己的“另一半”。
  \textbf{三个已标出的零点说明：原点必为零点，却未必唯一。}

  \medskip
  
\end{explanationbox}
\clearpage
\begin{explanationbox}
\textbf{\textcolor{CExplanation}{五、眼前一亮：把原点换成别的对称中心}}
\begin{center}\begin{tikzpicture}[font=\small,>=stealth,line join=round]
\draw[gray!40,rounded corners=3pt] (0,0) rectangle (14.2,8.5);\node at (7.1,7.9) {把原点换成 $(a,b)$：中心横坐标处只能取中心高度};
\begin{scope}[shift={(1.6,1.8)},x=1.4cm,y=0.8cm]\draw[->,gray] (-0.2,0)--(4.2,0) node[right]{$x$};\draw[->,gray] (0,-0.3)--(0,6.1) node[above]{$y$};\draw[blue,thick,domain=0.7:3.3,samples=150] plot(\x,{(\x-2)^3+3});\draw[gray,dashed] (2,0)--(2,5.7);\draw[gray,dotted] (0,3)--(3.9,3);\fill[orange!85!black] (2,3) circle(2.2pt) node[above left=5pt] {$S(2,3)$};\fill[orange!85!black] (1,2) circle(2.2pt) node[below right=5pt] {$(1,2)$};\fill[orange!85!black] (3,4) circle(2.2pt) node[above right=5pt] {$(3,4)$};\node[below=5pt] at (2,0){$2$};\node[left=5pt] at (0,3){$3$};\end{scope}\node[align=left,anchor=west] at (8.4,4.2){图取 $f(x)=(x-2)^3+3$\\[11pt]对称中心为 $(2,3)$\\[11pt]$f(2+t)+f(2-t)=6$\\[11pt]$x=2$ 有定义，所以 $f(2)=3$};
\node[text width=12.8cm,align=center] at (7.1,.65) {对称点 $(1,2)$、$(3,4)$ 的中点正是 $(2,3)$；一般情形仍要先确认 $x=a$ 有定义。};
\end{tikzpicture}\end{center}



  如果某个函数的图像不是关于原点，而是关于点 $(a,b)$ 中心对称，
  且它在 $x=a$ 处有定义，那么它在这条竖线上的点只能是 $(a,b)$，即
  \[
    \boxed{f(a)=b.}
  \]
  为什么？点 $(a,c)$ 关于 $(a,b)$ 对称后变为 $(a,2b-c)$，
  两个点的横坐标仍是 $a$；由于函数在同一横坐标处只能有一个值，
  必须有 $c=2b-c$，即 $c=b$。

  也可以用一个容易记的式子表示：
  \[
    f(a+t)+f(a-t)=2b.
  \]
  如果 $a$ 在定义域内，令 $t=0$，立刻得到 $2f(a)=2b$。
  \textbf{本节的 $f(0)=0$，就是取 $a=0,b=0$ 的特殊情形。}

  \medskip
  
\end{explanationbox}
\clearpage
\begin{explanationbox}
\textbf{\textcolor{CExplanation}{六、最后把最容易忘记的条件说清楚}}
\begin{center}\begin{tikzpicture}[font=\small,>=stealth,line join=round]
\draw[gray!40,rounded corners=3pt] (0,0) rectangle (14.2,8.5);\node at (7.1,7.9) {易错一：关于原点对称，不等于图像一定包含原点};
\node at (3.5,6.85){$y=x^3$：$0$ 可以代入};\node at (10.6,6.85){$y=1/x$：$0$ 不能代入};\begin{scope}[shift={(3.5,3.7)},x=1cm,y=1cm]\draw[->,gray] (-2.2,0)--(2.2,0) node[right]{$x$};\draw[->,gray] (0,-2.2)--(0,2.2) node[above]{$y$};\draw[blue,thick,domain=-1.25:1.25,samples=150] plot(\x,{(\x)^3});\fill[orange!85!black] (0,0) circle(2.2pt) node[below right=5pt] {$O$};\end{scope}\begin{scope}[shift={(10.5,3.7)},x=1cm,y=1cm]\draw[->,gray] (-2.2,0)--(2.2,0) node[right]{$x$};\draw[->,gray] (0,-2.2)--(0,2.2) node[above]{$y$};\draw[blue,thick,domain=-2:-0.5,samples=150] plot(\x,{1/\x});\draw[blue,thick,domain=0.5:2,samples=150] plot(\x,{1/\x});\fill[orange!85!black] (1,1) circle(2.2pt) node[above right=5pt] {$(1,1)$};\fill[orange!85!black] (-1,-1) circle(2.2pt) node[below left=5pt] {$(-1,-1)$};\node[below right=5pt] at (0,0){$O$};\end{scope}\node at (3.5,1.25){是奇函数，且 $f(0)=0$};\node at (10.6,1.25){是奇函数，但 $f(0)$ 不存在};
\node[text width=12.8cm,align=center] at (7.1,.65) {两图横纵同尺度；右图原点仅是坐标轴交点，不是函数图像上的点。};
\end{tikzpicture}\end{center}



  奇函数\textbf{不一定有} $f(0)$。
  例如 $f(x)=\dfrac1x$，它满足
  \[
    f(-x)=\frac1{-x}=-\frac1x=-f(x)\qquad(x\ne0),
  \]
  确实是奇函数，但 $x=0$ 时分母为零，$f(0)$ 根本不存在。
  所以正确的说法永远是：\textbf{奇函数若在 $0$ 处有定义，才必有 $f(0)=0$。}
\end{explanationbox}
